\documentclass[11pt,a4paper]{article}
\usepackage[etiket=en]{gtu-el-yazisi}
\begin{document}
\notbilgi{Kontrol Sistemleri}{Chapter 7 Examples: Steady-State Error Problems}{Chapter-7-Examples.pdf}
\textit{Not: Orijinal dosyadaki basılı sayfalar (s. 7) dijitalleştirilmedi; yalnızca el yazısı sayfalar aktarıldı.}

\sayfa{1}
\sekilnotu[1]{Sayfanın üst kısmında basılı soru 6 bulunuyor; dijitalleştirilmedi (Tip 3 birim geri beslemeli sisteme $25t^{3}u(t)$ girişi uygulanıyor; $G(s)=\dfrac{210(s+4)(s+6)(s+11)(s+13)}{s^{3}(s+7)(s-14)(s+19)}$; konumdaki kalıcı hal hatası isteniyor; Bölüm 7.3).}
\[
t^{n}\donusum\frac{n!}{s^{n+1}}\qquad\qquad E(s)=\frac{R(s)}{1+G(s)}
\]
\begin{align*}
25\cdot t^{3}&\donusum\frac{6}{s^{4}}\cdot25=\frac{150}{s^{4}}\\[4pt]
\lim_{s\to0}sE(s)&=\frac{\dfrac{150}{{\not s}^{3}}}{\dfrac{210(s+4)(s+6)(s+11)(s+13)}{{\not s}^{3}(s+7)(s-14)(s+19)}}\\[8pt]
&=\frac{150}{\dfrac{(210)(4)(6)(11)(13)}{(7)(14)(19)}}\\[8pt]
&=0{,}38
\end{align*}
\sekilnotu[1]{Sayfanın ortasında basılı soru 12 bulunuyor; dijitalleştirilmedi (birim geri beslemeli sistemde $G(s)=\dfrac{K(s+2)(s+4)(s+6)}{s^{2}(s+5)(s+7)}$; statik hata sabiti $10\,000$ olacak şekilde $K$ değeri isteniyor; Bölüm 7.4).}
\begin{align*}
E(s)&=\frac{R(s)}{1+G(s)}\\
R(s)&=\frac{1}{s}\\
\Rightarrow e_{ss}&=\lim_{s\to0}\frac{1}{\underbrace{1+G(s)}_{k_p}}\\
k_p&=\lim_{s\to0}G(s)\\[4pt]
k_a&=\lim_{s\to0}s^{2}G(s)=\lim_{s\to0}\frac{k(s+2)(s+4)(s+6)}{{\not s}^{2}(s+5)(s+7)}\\[4pt]
10\,000&=\frac{k\cdot2\cdot4\cdot6}{5\cdot7}\\
\Rightarrow\ 7291&=k
\end{align*}

\sayfa{2}
\sekilnotu[2]{Sayfanın üst kısmında basılı soru 13 ve Şekil P7.4 bulunuyor; dijitalleştirilmedi (ileri yolda $\dfrac{5}{s(s+1)(s+2)}$ bloğu, geri besleme yolunda $(s+3)$ bloğu; a. $K_p$, $K_v$, $K_a$ bulunacak, b. $50u(t)$, $50tu(t)$ ve $50t^{2}u(t)$ girişleri için kalıcı hal hatası, c. sistem tipi; Bölüm 7.3).}
\begin{align*}
\frac{\dfrac{5}{s(s+1)(s+2)}}{\underbrace{1+\dfrac{5(s+3)}{s(s+1)(s+2)}}_{\dfrac{s(s+1)(s+2)+5(s+3)}{s(s+1)(s+2)}}}
&=\frac{5}{s\underbrace{(s+1)(s+2)}_{s^{2}+3s+2}+5(s+3)}\\[10pt]
&=\frac{5}{s^{3}+3s^{2}+2s+5s+15}=\frac{5}{s^{3}+3s^{2}+7s+15}=G(s)
\end{align*}
\elcizimi[0.6]{sekiller/final-calismasi-Chapter-7-Examples-s2-1.png}[2]
\begin{align*}
k_p&=\lim_{s\to0}G(s)=\frac{1}{3}\\
k_v&=0\\
k_a&=0
\end{align*}
for

b) $E_{ss}=\dfrac{50}{1+k_p}=\dfrac{50}{1+\frac{1}{3}}=37{,}5$

for $50\,tu(t)$
\[
E_{ss}=\lim_{s\to0}\frac{50}{s}=\infty
\]
for $50\,t^{2}u(t)\;=\;E_{ss}=\infty$

\sayfa{3}
\sekilnotu[3]{Sayfanın üst kısmında basılı soru 15 ve Şekil P7.5 bulunuyor; dijitalleştirilmedi (ileri yolda $\dfrac{100(s+2)}{s(s+5)}$ ve $\dfrac{1000}{s}$ blokları, ilk blok etrafında 10 kazançlı iç geri besleme ve birim dış geri besleme; sistem tipi isteniyor; Bölüm 7.3). Elle eklenen: ilk bloğu ve 10 bloğunu çevreleyen büyük bir kalem çizgisi ve bu çizgiden ``$G(s)=$'' ifadesine giden ok.}
\[
G(s)=\frac{\dfrac{100(s+2)}{s(s+5)}}{1+\dfrac{1000(s+2)}{s(s+5)}}=\frac{100(s+2)}{s^{2}+\underbrace{5s+1000s}_{1005\,s}+2000}
\]
\[
k_v=\lim_{s\to0}\frac{100(s+2)\,(1000)}{s^{2}+1000s\footnotemark+2000}
\]
\footnotetext{Claude notu: muhtemelen $1005s$ olmalı.}
\sekilnotu[3]{Sayfanın ortasında basılı soru 16 ve Şekil P7.6 bulunuyor; dijitalleştirilmedi ($G_1(s)$ sırasıyla Tip 0, Tip 1 ve Tip 2 iken basamak girişte sıfır kalıcı hal hatası için ileri besleme $G_2(s)$ üzerindeki kısıtlar soruluyor; plant $\dfrac{20}{s(s+3)}$; Bölüm 7.3).}
\elcizimi[0.9]{sekiller/final-calismasi-Chapter-7-Examples-s3-1.png}[3]
\[
G(s)=\underbrace{\Bigl(\underbrace{G_1}_{\frac{1}{s+3}}+R\,\underset{\searrow\,\frac{1}{s}}{G_2}\Bigr)\Bigl(\frac{20}{s(s+3)}\Bigr)}_{\text{Type 1}}
\]

\sayfa{4}
\sekilnotu[4]{Sayfanın üst kısmında basılı soru 22 ve Şekil P7.1 bulunuyor; dijitalleştirilmedi (birim geri beslemeli sistemde $G(s)=\dfrac{K(s^{2}+3s+30)}{s^{n}(s+5)}$; $10tu(t)$ girişi ile çıkış arasında kalıcı halde $1/6000$ hata isteniyor; a. $K$ ve $n$ bulunacak, b. $K_p$, $K_v$, $K_a$ nedir; Bölüm 7.4).}
\[
\lim_{s\to0}\frac{10}{s+\underbrace{sG(s)}_{k_v}}=\frac{1}{6000}
\]
$n=1$
\[
\lim_{s\to0}\;=\;\frac{sk(s^{2}+3s+30)}{{\not s}(s+5)}\;=\;\frac{30k}{5}
\]
\[
\lim_{s\to0}\frac{10}{s+\underbrace{\dfrac{30k}{5}}_{6}}=\frac{1}{6000}
\]
\begin{align*}
k_p&=\infty&&\qquad\Rightarrow\ k=10\,000\\
k_v&=6k=60\,000\\
k_a&=0
\end{align*}

\sayfa{5}
\sekilnotu[5]{Sayfanın üst kısmında basılı soru 27 bulunuyor; dijitalleştirilmedi (birim geri beslemeli sistemde $G(s)=\dfrac{K(s+\alpha)}{(s+\beta)^{2}}$; birim basamak girişte kalıcı hal hatası $0{,}1$, sönüm oranı $0{,}5$, doğal frekans $\sqrt{10}$ olacak şekilde $K$, $\alpha$ ve $\beta$ isteniyor; Bölüm 7.4).}
\begin{gather*}
\frac{\omega_n^{2}}{s^{2}+2\zeta\omega_n s+\omega_n^{2}}\\[4pt]
\frac{10}{s^{2}+\sqrt{10}\,s+10}
\end{gather*}
\begin{align*}
e(\infty)&=\frac{1}{1+k_p}=0{,}1\\
k_p&=\lim_{s\to0}G(s)=\frac{k\alpha}{\beta^{2}}\\
\zeta&=0{,}5\qquad\omega_n=\sqrt{10}\\
10&=1+k_p\qquad k_p=9\qquad\frac{k\alpha}{\beta^{2}}=9
\end{align*}
\sekilnotu[5]{Sayfanın ortasında basılı soru 32 bulunuyor; dijitalleştirilmedi (birim geri beslemeli sistemde $G(s)=\dfrac{K}{s(s+a)}$; a. $K_v=1000$ ve yüzde 20 aşım için $K$ ve $a$, b. kalıcı halde yüzde 1 hata ve yüzde 10 aşım için $K$ ve $a$; Bölüm 7.4).}
\begin{align*}
&=\frac{\dfrac{k}{s(s+a)}}{1+\dfrac{k}{s(s+a)}}\\[8pt]
&=\frac{k}{s(s+a)+k}=\frac{\overset{\omega_n^{2}}{k}}{s^{2}+as+\underbrace{k}_{\omega_n^{2}}}
\end{align*}
\begin{gather*}
2\zeta\omega_n=a\qquad\omega_n=\sqrt{k}\\
2\zeta\sqrt{k}=a\\
\zeta=\frac{a}{2\sqrt{k}}
\end{gather*}
\begin{align*}
M_p&=e^{-t_p\sigma}\\
&=e^{-\frac{\pi}{\omega_d}\sigma}
\end{align*}
\begin{gather*}
\sigma=\zeta\omega_n\\
t_p=\frac{\pi}{\omega_d}\\
\omega_d=\omega_n\sqrt{1-\zeta^{2}}
\end{gather*}
\[ k_v=\lim_{s\to0}sG(s)=\frac{sk}{s(s+a)}=\frac{k}{a}=1000 \]
\begin{gather*}
M_p=e^{-\frac{\pi\zeta}{\sqrt{1-\zeta^{2}}}}=0{,}2\qquad\qquad\frac{\omega_n^{2}}{2\zeta\omega_n}=1000\\
\zeta=0{,}45\qquad\qquad\omega_n=2000\zeta\\
\omega_n=2000\zeta=912\qquad\qquad k=\omega_n^{2}=831\,744\\
\frac{k}{a}=1000\ \Rightarrow\ a=831{,}744
\end{gather*}
\sekilnotu[5]{Sayfanın alt kısmında basılı soru 39 ve Şekil P7.15 bulunuyor; dijitalleştirilmedi (blok diyagram: ileri yolda $\dfrac{K_1(s+2)}{s+3}$ ve $\dfrac{K_2}{s(s+4)}$ blokları, $D(s)$ bozucusu ikinci toplama noktasına giriyor, birim geri besleme; birim basamak bozucudan kaynaklanan kalıcı hal hatası bileşeni $-0{,}00001$, birim rampa girişten kaynaklanan bileşen $0{,}002$ olacak şekilde $K_1$ ve $K_2$ isteniyor; Bölüm 7.5).}
\begin{align*}
e_D(\infty)&=-0{,}000012\footnotemark\\
&=-\frac{3}{2k_1}\\
\Rightarrow\ k_1&=125\,000
\end{align*}
\footnotetext{Claude notu: soruda $-0{,}00001$ verilmiş; muhtemelen $-0{,}00001$, $k_1=150\,000$ ve $k_2=0{,}02$ olmalı.}
\begin{align*}
\frac{1}{2k_1}&=0{,}00001\\
k_1&=\frac{\text{\okunamadi{1}}}{2\cdot0{,}00001}\\
&=\text{\okunamadi{50\,000}}
\end{align*}
\begin{align*}
e_D(\infty)&=-\lim_{s\to0}\frac{sG_2D}{1+G_1G_2}\\[6pt]
&=-\lim_{s\to0}\frac{s\left(\dfrac{k_2}{s(s+4)}\right)\dfrac{1}{s}}{1+\dfrac{k_1(s+2)}{s+3}\cdot\dfrac{k_2}{s(s+4)}}\\[10pt]
&=\frac{k_2\,\text{\okunamadi{}}}{s(s+3)(s+4)+k_1(s+2)k_2}\\[8pt]
&=\frac{{\not k}_{2}\underbrace{(s+3)}_{3}}{2k_1{\not k}_{2}}=-\frac{3}{2k_1}
\end{align*}
\formulkutu{e_R(\infty)=0{,}002=\frac{6}{k_1k_2}\Rightarrow k_2=0{,}024}

\sayfa{6}
\sekilnotu[6]{Sayfanın üst kısmında basılı soru 43 ve Şekil P7.18 (Sistem 1 ve Sistem 2) bulunuyor; dijitalleştirilmedi (her sistem için uygun statik hata sabiti ile birim basamak, rampa ve parabolik girişler için kalıcı hal hatası $r(\infty)-c(\infty)$ isteniyor; Bölüm 7.5). Elle eklenen: Sistem 1 şeklinin üzerinde kalemle bir çarpı işareti ve $10$ bloğu çevresinde birkaç çizgi var.}
\[
G(s)=\frac{5(s+4)}{(s+3)(s+7)}
\]
\sekilnotu[6]{Sayfanın ortasında basılı soru 45 (Şekil P7.20) ve soru 47 (Şekil P7.22, DVR modeli) bulunuyor; dijitalleştirilmedi.}
\sekilnotu[6]{Sayfanın altında basılı soru 51 bulunuyor; dijitalleştirilmedi (Şekil P7.13: a. plant değişimlerine duyarlılığın $S_{T:P}=\dfrac{P}{T}\dfrac{\delta T}{\delta P}=\dfrac{1}{1+L(s)}$ olduğu gösterilecek, burada $L(s)=G(s)P(s)H(s)$ ve $T(s)=\dfrac{C(s)}{R(s)}=\dfrac{F(s)}{H(s)}\dfrac{L(s)}{1+L(s)}$; b. $S_{T:P}(s)+\dfrac{T(s)H(s)}{F(s)}=1$ olduğu gösterilecek).}
\[
T=\frac{F}{H}\,\frac{L}{1+L}
\]
\formulkutu{T=\frac{F}{{\not H}}\,\frac{GP{\not H}}{1+GPH}}
\[
T=\frac{FGP}{1+GPH}
\]
\[
\frac{P}{T}=\frac{P}{\dfrac{F(s)}{H(s)}\,\dfrac{L(s)}{1+L(s)}}
\]

\end{document}
